\documentclass{article}
\usepackage{tikz}
\usetikzlibrary{positioning, arrows.meta, shapes.geometric, backgrounds}
\usepackage{lscape}
\usepackage{geometry}

% \usepackage[margin=0.5cm]{geometry}

\begin{document}

\begin{landscape}
\begin{figure}[h!]
\centering
\begin{tikzpicture}[
    node distance=0.8cm,
    inputlayer/.style={rectangle, draw=black!60, fill=green!20, rounded corners=0.3cm, 
                      minimum width=1.5cm, minimum height=2.5cm, align=center, font=\large, thick},
    linearlayer/.style={rectangle, draw=black!60, fill=blue!20, rounded corners=0.3cm, 
                       minimum width=1.5cm, minimum height=2.5cm, align=center, font=\normalsize, thick, font=\large},
    activationlayer/.style={rectangle, draw=black!60, fill=red!20, rounded corners=0.3cm, 
                           minimum width=1.5cm, minimum height=2.5cm, align=center, font=\normalsize, thick,font=\large},
    outputlayer/.style={rectangle, draw=black!60, fill=purple!20, rounded corners=0.3cm, 
                       minimum width=1.5cm, minimum height=2.5cm, align=center, font=\normalsize, thick,font=\large},
    arrow/.style={-Stealth[scale=1.1], thick, black!80},
]

% Layers
\node[inputlayer] (input) {Input\\(B, N, F)};
\node[linearlayer, right=of input] (linear1) {Linear\\$F \rightarrow 256$};
\node[activationlayer, right=of linear1] (relu1) {ReLU};

\node[linearlayer, right=of relu1] (linear2) {Linear\\$256 \rightarrow 128$};
\node[activationlayer, right=of linear2] (relu2) {ReLU};

\node[linearlayer, right=of relu2] (linear3) {Linear\\$128 \rightarrow K \times 3$};

\node[outputlayer, right=of linear3] (output) {Output\\$(B, N, K, 3)$\\ reshape};

% Arrows
\draw[arrow] (input) -- (linear1);
\draw[arrow] (linear1) -- (relu1);
\draw[arrow] (relu1) -- (linear2);
\draw[arrow] (linear2) -- (relu2);
\draw[arrow] (relu2) -- (linear3);
\draw[arrow] (linear3) -- (output);

% Background
\begin{scope}[on background layer]
    \draw[fill=gray!5, rounded corners=0.5cm] 
        ([xshift=-1cm,yshift=1cm]input.north west) rectangle 
        ([xshift=1cm,yshift=-1cm]output.south east);
\end{scope}

\end{tikzpicture}
\end{figure}
\end{landscape}

\end{document}